\chapter{学习动力学的重建：从变分目标到受约束随机更新}
\label{chp:dynamical_reconstruction}

前章已经把计算底空间拆成对象身份、任务度量与方向接口，本章进一步追问这些结构怎样在经验到来时改变。我们不把深度学习训练直接说成某种物理场方程的实现，也不由损失下降推出智能任务成功；我们要建立的是一组可比较、可离散和可反驳的学习动力学。为此，我们分别定义概率模型、参数状态、结构状态、数据分布、任务目标、更新时钟与资源账本，再说明变分自由能、信息几何、随机梯度和耗散语言各自能够支持什么结论。HSF-HD在这里给出一般状态—结构耦合的模型族，AgentOS中的$h(t)$、$E$、$\Theta$、$\Psi$及外部记忆$M_e$只是其中一种工程实例。只有在对象、单位与假设明确时，几何和热力学类比才用于解释；它们不能代替收敛证明、泛化检验或实际功耗测量。

\section{从变分目标到状态—结构耦合，而非严格同构}
\label{sec:free_energy_hsf_hd}

我们先固定对象。设观测$x\in\mathcal X$、目标$y\in\mathcal Y$，训练时刻为$k\in\mathbb N$，模型参数$\theta_k\in\Theta\subseteq\mathbb R^d$，预测分布为$p_{\theta_k}(y\mid x)$。数据批次$B_k$由可能随时间变化的分布$q_k(x,y)$产生。若系统还维护可修改的局部结构网络，则记其结构状态为$G_k\in\mathcal G$；若维护全局分布表示，则记快状态为$h_k\in\mathcal H$。$G_k$描述可寻址节点、边、权限和接口，$h_k$描述分布式激活；二者通过任务相关身份映射、绑定算子和读出接口结合，不能因为都能影响输出就视为同一个对象。外部记忆$M_e$、边界控制与卸载记录属于可选的AgentOS实现，不是所有智能系统都必须具有的本体部件。

对给定批次，我们定义无量纲的经验目标
\begin{equation}
J_k(\theta,G)=\frac{1}{|B_k|}\sum_{(x,y)\in B_k}\ell\!\left(p_{\theta,G}(\cdot\mid x),y\right)+\lambda_k\Omega(\theta,G)+C_k(\theta,G),
\end{equation}
其中$\ell$是预测损失，$\Omega$是显式复杂度惩罚，$C_k$是任务约束的软罚项；若三者原有单位不同，就必须先用各自基准量归一化。这个$J_k$是工程定义的训练目标，不是焦耳能量，也不是香农熵。若另有潜变量$z$和近似后验$r_\phi(z\mid x)$，则可以定义变分自由能
\begin{equation}
\mathcal F_k(\theta,\phi)=\mathbb E_{q_k(x)}\!\left[D_{\mathrm{KL}}\!\left(r_\phi(z\mid x)\Vert p_\theta(z)\right)-\mathbb E_{r_\phi}\log p_\theta(x\mid z)\right].
\end{equation}
在密度存在且期望有限的条件下，它是负对数证据的上界。复杂度项和准确度项的分解是一个概率推断结果；只有所讨论模型确实包含这些分布时才成立。普通监督损失可以与它共享优化形式，却不因此与它严格同构，更不能从数值相近推出机制相同。

HSF-HD把学习写成状态—结构耦合更新族
\begin{equation}
(h_{k+1},\theta_{k+1},G_{k+1})=\mathcal U_{\Delta\tau_k}
\left(h_k,\theta_k,G_k,u_k,q_k,m_k\right),
\end{equation}
其中$u_k$是外部输入，$m_k$是度量与资源记录，$\Delta\tau_k>0$是离散步长。该式首先是一项定义。进一步假设$G_k$在若干快步内固定、$h_k$先弛豫，而$\theta_k$与$G_k$在慢时钟更新，才可使用双时间尺度分析。若$q_k$、目标权重$\lambda_k$或约束集合随时间改变，则$J_k$本身也在改变，不能把$J_{k+1}(\theta_{k+1})<J_k(\theta_k)$简单解释为同一目标上的下降。

我们因此把历史作用量语言降为候选建模接口：可以构造轨迹泛函$\mathcal A[\xi]=\sum_k L_k(\xi_k,\xi_{k+1})\Delta\tau_k$来比较路径，其中$\xi_k=(h_k,\theta_k,G_k)$；但其驻点仅在给定边界、可微性和变分集合下成立。真实训练通常含随机批次、投影、截断和图编辑，未必来自单一光滑作用量。可检验的主张应当是：在声明的任务分布与预算下，某种耦合更新比对照更新取得更低校准误差、更少约束违规或更好的样本外性能。反例也清楚：当数据标签系统性错误时，$J_k$持续下降仍可增加任务失败；当复杂度惩罚过强时，变分上界改善也可能以欠拟合为代价。由此，目标下降是局部优化证据，而不是智能成功或物理同构的证明。

这一区分还决定怎样解释“准确度”和“复杂度”。前者可以是对数似然的期望，也可以是具体任务的正确率，但两者不是同一量；后者可以是KL正则、参数范数、结构边数或描述长度，也不能彼此替代。若研究者改变了损失定义、样本权重或结构搜索空间，比较基准随之改变，导数中必须出现这些显式变化。我们要求每次运行保存目标版本、数据快照、随机种子和约束配置，并在测试集之外保留压力输入：标签偏移检验预测项，结构扰动检验$G$的恢复能力，预算收紧检验控制器是否只是用更多资源换取表面改进。这样，HSF-HD的作用不是为已有算法赋予宏大名称，而是强迫模型把状态、结构、目标、输入和证据链放在同一个可审计时间轴上。

这种账本也让同一公式在不同任务中的含义保持可追踪，而不是依靠术语相似性迁移结论。

\section{信息几何只度量分布可辨识性}
\label{sec:cognitive_einstein_information_geometry}

我们把参数空间中的几何限制在它能够严格表达的范围内。设$\theta$位于开集$\Theta\subseteq\mathbb R^d$，模型族$p_\theta(o)$对$\theta$二次可微，概率密度具有共同支撑，且微分与积分可以交换。得分向量$s_\theta(o)=\nabla_\theta\log p_\theta(o)$无量纲，费希尔信息矩阵定义为
\begin{equation}
F(\theta)=\mathbb E_{o\sim p_\theta}\left[s_\theta(o)s_\theta(o)^{\mathsf T}\right]\in\mathbb R^{d\times d}.
\end{equation}
它是半正定矩阵；当参数局部可辨识时才为正定。若参数带单位，$F_{ij}$的量纲是$\theta_i^{-1}\theta_j^{-1}$，因此二次型$\delta\theta^{\mathsf T}F\delta\theta$无量纲。对足够小的$\delta\theta$，KL散度满足
\begin{equation}
D_{\mathrm{KL}}(p_\theta\Vert p_{\theta+\delta\theta})
=\frac12\delta\theta^{\mathsf T}F(\theta)\delta\theta+o(\|\delta\theta\|^2).
\end{equation}
这说明$F$度量的是邻近预测分布的可辨识性，而不是知识网络的全部拓扑，也不是实际材料的时空度量。参数对称、冗余神经元或尺度不变性会使$F$奇异；此时应使用阻尼、商空间或广义逆，不能假定逆矩阵天然存在。

给定可微目标$J(\theta)$，我们在局部KL步长约束$\frac12\delta\theta^{\mathsf T}F\delta\theta\leq\varepsilon$下最小化一阶近似$\nabla J^{\mathsf T}\delta\theta$。拉格朗日乘子法给出
\begin{equation}
\delta\theta=-\alpha F(\theta)^{-1}\nabla J(\theta),
\end{equation}
其中$\alpha$由约束半径决定。这是自然梯度的连续局部推导。若使用$F_\lambda=F+\lambda I$，则方向对应带阻尼的工程设计；若只保留对角线或分块近似，则参数化不变性只部分保留。所谓“最快”也必须附带度量：它是在局部预测分布距离下的一阶最陡下降，不是墙钟时间、浮点操作数或实际能耗意义上的最快。

结构状态$G$不能直接塞进这个连续参数流形。边的增删、工具权限变化和记忆对象合并属于离散编辑，适合写成候选操作集$\mathcal E(G)$上的受约束选择；分布状态$h$可以在固定结构下连续演化。我们把联合更新拆成三步：先在$G_k$固定时计算$h$与$\theta$的候选变化，再由结构控制器根据任务收益、回滚能力和边界约束选择编辑$e_k\in\mathcal E(G_k)$，最后通过绑定表重新建立局部身份与全局表示之间的读出关系。这样的分层避免从脑区或几何类比推出机制等同，也使结构失败能够单独回滚。

历史上的“状态—结构耦合方程”可以在本章中用一个可验证模型替代。设慢结构参数为$\zeta\in\mathbb R^r$，正定任务度量为$M(\zeta)$，则
\begin{align}
\dot\theta&=-M(\zeta)^{-1}\nabla_\theta J(\theta,\zeta;q_\tau),\\
\dot\zeta&=-K(\zeta)\nabla_\zeta\bigl[J(\theta,\zeta;q_\tau)+R(\zeta)\bigr],
\end{align}
其中$K$也是正定迁移率，$q_\tau$允许环境漂移。若$q_\tau$固定且$M,K$无显式时间变化，则目标沿轨迹的导数为两个非正二次型之和；若$q_\tau$变化，还必须加上$\partial J/\partial\tau$。因此只能在附加条件下得到单调性。即使达到驻点，也只能说明一阶必要条件，不能断言全局最优、泛化最佳或未来输入的耗散最小。验证应比较完整费希尔、对角近似、经验梯度平方和欧氏基线，报告条件数、信赖域违反率、训练损失与样本外指标；若几何近似没有提供可重复改进，几何解释就不应被当成机制事实。

度量本身也可能成为学习对象，这时必须把它的变化计入分析。若$M=M(\theta,\tau)$，则相邻两步的“距离”使用了不同标尺；把两次二次型直接相减没有固定意义。我们可以冻结一个审计窗口内的参考度量，或使用把度量演化显式写入的协变导数，但不能悄悄沿用静态证明。工程上更稳妥的做法是给度量快照编号，记录其估计样本和有效期，并在分布漂移超过阈值时失效。一个尤其重要的反例是重参数化：两个网络产生完全相同的预测分布，却可能具有不同的欧氏梯度和参数范数；费希尔几何能解释局部预测不变性，却不能替代结构身份、因果角色或记忆来源的判定。因此，本节只把它当成分布读出层的局部标尺。

\section{从随机梯度到有条件的朗之万模型}
\label{sec:statistical_physics_view}

随机梯度首先是数据抽样算法，而不是自动成立的热噪声模型。设固定目标$J(\theta)=\mathbb E_{\xi\sim q}[\ell(\theta;\xi)]$，小批量无偏估计为$\widehat g_k=\nabla J(\theta_k)+\epsilon_k$，并假设$\mathbb E[\epsilon_k\mid\theta_k]=0$、条件协方差为$\Sigma(\theta_k)/b$，其中$b$为批量大小。基本更新
\begin{equation}
\theta_{k+1}=\theta_k-\eta_k P_k\widehat g_k
\end{equation}
使用正定预条件器$P_k$。在梯度利普希茨、二阶矩有界、步长足够小等条件下，可以得到期望下降界；非凸情形通常只保证某种平均梯度范数趋小，不保证抵达全局极小值。动量、数据重用、裁剪和自适应预条件会改变噪声的相关性，因而不能不经检验就把SGD等同于平衡朗之万过程。

若我们有意构造采样器，才定义连续随机微分方程。令目标密度$\pi(\theta)\propto\exp[-J(\theta)/T]$，$T>0$是与$J$同单位的算法尺度；当$J$已经无量纲时$T$也无量纲。对常数正定矩阵$P$，过阻尼模型为
\begin{equation}
d\theta_\tau=-P\nabla J(\theta_\tau)d\tau+\sqrt{2TP}\,dW_\tau.
\end{equation}
在边界通量为零、归一化常数有限等条件下，$\pi$是其平稳分布。若$P=P(\theta)$随状态改变，伊藤形式还需要由$\nabla\!\cdot P$产生的漂移修正；忽略它一般会改变平稳分布。这里的$T$控制探索方差，不是芯片温度、神经活动强度或焦耳能耗。实际训练是否近似此方程，要通过增量分布、噪声协方差、时间相关和稳态样本检验，而不能靠“有随机性”这一表面相似来决定。

采用Euler--Maruyama离散化，固定步长$\Delta\tau$时有
\begin{equation}
\theta_{k+1}=\theta_k-\Delta\tau P_k\nabla J_k(\theta_k)
+\sqrt{2T_k\Delta\tau}\,P_k^{1/2}z_k,
\qquad z_k\sim\mathcal N(0,I).
\end{equation}
若还使用小批量梯度，抽样噪声与显式噪声必须分账；前者通常依赖$\theta$且非高斯，不能简单并入同一个温度。稳定性要求步长相对于预条件后的局部曲率足够小，约束参数还需投影或反射边界。探索可以帮助越过局部势垒，却会增加估计方差；温度退火可能改善最终集中，也可能过早冻结。我们因此把$T_k$、批量$b_k$、步长$\eta_k$和预条件器版本都写入运行记录，使一次失败能够定位到具体控制量。

在AgentOS示例中，$h(t)$可作为快状态，$\Theta$作为可调结构参数，$E$必须注明究竟是无量纲预算、推理代价还是测得电能；$\Psi$若表示任务分布状态，也须给出状态空间和归一化规则。SPC可以管理探索日程，TCE可以执行受约束候选更新，CCM可以检查跨模块一致性，Boundary负责权限与可逆性，Offloading把长程证据写入$M_e$。这些模块可实现上述控制回路，但理论只要求存在相应功能，不要求所有智能系统采用这些名称或相同分解。

一个可证伪的实验至少包含四组：欧氏SGD、固定预条件器、状态相关自然梯度近似、显式朗之万更新。我们报告同等计算预算下的训练目标、校准误差、分布外风险、约束违规、参数位移和墙钟时间；焦耳能耗则由设备功率采样积分获得。若加噪只降低训练损失却恶化任务成功率，它不是“创造力提高”；若自然梯度减少步数却增加总能耗，也不能称为热力学更优。统计物理语言在这里提供候选随机模型，结论仍由数据和适用条件决定。

还要防止把随机性本身浪漫化。噪声可能帮助越过窄势垒，也可能破坏稀有但正确的表示；批次噪声可能提供正则效应，也可能只是数据排序造成的偏差。我们在固定预算下扫描温度与步长的联合网格，分别报告到达时间、稳态波动和最坏分位风险，并用关闭显式噪声、打乱批次相关、保持总方差三种消融区分来源。若任务目标在运行中变化，则比较动态遗憾而非静态终值。只有这些观测与随机微分模型的预测相符时，“朗之万近似”才获得经验支持；否则它只是方便的算法构造。

\section{受约束二阶近似的离散优化器}
\label{sec:engineering_mapping_optimizer}

我们把历史上的Sophia--HSF描述重建为一个明确的工程实例，而不宣称它是一般智能的必然更新律。设$m_k\in\mathbb R^d$为梯度一阶矩估计，$v_k\in\mathbb R_+^d$为非负曲率代理，$\varepsilon>0$为数值下界，$\gamma>0$为曲率缩放，$\rho>0$为逐坐标信赖域半径。对批次损失$J_k$，先计算$g_k=\nabla J_k(\theta_k)$并更新
\begin{align}
m_k&=\beta_1m_{k-1}+(1-\beta_1)g_k,\\
v_k&=\beta_2v_{k-1}+(1-\beta_2)\widehat c_k,
\end{align}
其中$\widehat c_k\geq0$可以是广义高斯—牛顿对角、经验费希尔对角或其他声明过偏差的估计器。负对数似然且标签按模型分布重采样时，得分平方的期望给出模型费希尔对角；这不等于任意损失的真实Hessian，也不保证与训练数据分布下的经验费希尔相同。估计频率、重采样规则和指数平均的滞后都必须纳入误差分析。

确定性候选步定义为
\begin{equation}
d_{k,i}=\operatorname{clip}\!\left(
\frac{m_{k,i}}{\max(\gamma v_{k,i},\varepsilon)},-\rho,\rho\right),
\qquad
\theta_{k+1}^{\mathrm{cand}}=\theta_k-\eta_k d_k.
\end{equation}
这里的截断是逐坐标信赖域设计：它直接给出$\|\theta_{k+1}^{\mathrm{cand}}-\theta_k\|_\infty\leq\eta_k\rho$，但不自动保证损失下降，更不证明参数映射的全局李普希茨性。要获得下降结论，需要目标在候选邻域内光滑、曲率代理误差有界，并用接受测试比较预测下降与实际下降。若测试失败，控制器缩小$\eta_k$或$\rho$并重算；若连续失败，则回滚到最近检查点。所谓“保护拓扑”在这里被落实为具体不变量，例如结构身份不变、边界权限不越界和参数步长有界，而不是一个未定义的物理性质。

若任务允许随机探索，我们另行生成$z_k\sim\mathcal N(0,I)$，定义噪声候选
\begin{equation}
n_{k,i}=\sigma_{k,i}z_{k,i},\qquad
\sigma_{k,i}^2=\frac{2T_k\eta_k}{\max(\gamma v_{k,i},\varepsilon)},
\end{equation}
并在独立安全半径$\rho_n$内截断。这个选择使低曲率坐标获得较大探索方差，但它只是预条件采样器的近似；状态相关方差、截断与接受规则都会改变其平稳分布。如果目标是优化而非精确采样，我们可以把它作为启发式扰动并通过消融验证；如果目标是从指定分布采样，则必须补上相应漂移修正或Metropolis接受步骤。确定性步和随机步分别记录，避免把性能变化错误归因于“温度”。

结构更新不能与参数更新悄然混合。对候选结构编辑$e\in\mathcal E(G_k)$，TCE先在隔离状态上执行，CCM检查绑定、读出和依赖一致性，Boundary验证权限、资源上限与回滚路径；只有通过两阶段提交，$G_{k+1}$才取候选值。连续参数候选则以版本号关联到同一结构快照。若并发训练期间$G$或数据分布发生变化，旧曲率代理$v_k$可能失效，系统必须降权或重置，不能继续把它当成当前度量。外部记忆$M_e$保存估计器版本、批次摘要、随机种子和接受结果，使失败复现实验成为可能。

实现验收分为数值、安全、任务和资源四本账。数值账检查有限值、分母下界、步长界及回滚一致性；安全账检查权限、不可修改对象和结构不变量；任务账比较训练目标、验证风险、校准与分布外表现；资源账记录显存、吞吐、延迟、设备功率积分和通信量。算法代价可用浮点操作数或墙钟时间表示，实际能耗以焦耳测得，两者不能互换。对角近似在强耦合峡谷中可能方向错误，逐坐标裁剪也可能破坏旋转不变性；显式噪声在小数据任务上可能压倒信号。这些都是设计边界。只有在预注册预算和多随机种子对照中，受约束二阶近似稳定提高目标指标时，我们才说它是有效的AgentOS优化组件，而不把它提升为所有学习系统的唯一动力学。

最后，优化器的输出只是候选参数版本。部署前仍需独立的任务验收与边界验收；训练目标的下降不能替代端到端成功率，也不能授权系统扩大工具权限。这个分离使学习失败、结构失败和执行失败能够分别归因。
